\section{Formal Verification}
\label{ch:verify}

\subsection{Object and Method of Verification}

The algebraic skeleton of the argument chain has been machine-verified with Lean 4 and Mathlib, fully implemented with zero \texttt{sorry} and zero warnings; Lean 4 has become an active platform at the intersection of machine-assisted proof and machine learning \citep{yang2023}. The verification covers Definitions 1 to 4, Lemmas 1 to 3, Proposition 4 and Theorem 5, and the final theorem \texttt{finale}, which composes the generalization side and the carrier side into a single conjunction. Definition~2 is formalized in both of its readings rather than only in the scalar one: the toy carrier carries the indexed family itself, the constant family that Lemma~1 excludes, the frozen family with one entry per position, and the first-order algebra of the general form. Two claims about the index are therefore machine-checked rather than prose: that a frozen family is not a loophole for the algebra, since on a single configuration it satisfies every stationarity condition by construction while no scalar family does; and, sharper, that one position whose demand ratio changes between two admissible configurations excludes every frozen family, whatever the number of entries, which is the machine-checked form of the sentence that one scalar is a parameter table and a million parameters are a parameter table as well. On the toy carrier the general form's premise is derived rather than assumed: differentiating the objective along the coordinate line of one position gives the per-position coefficient \eqref{eq:focgeneral} of an arbitrary even differentiable potential, and a local minimum along that line satisfies the resulting first-order condition, which is \eqref{eq:foc} with the normalization of constant factors made explicit. Outside the machine guarantee remain the continuous-side correspondence that turns the variational objective into pointwise form and the operator-valued couplings that the general form does not claim; Section~\ref{sec:layers} records them as Tier~C. The formalization adopts a toy carrier --- four positions, structure fields as functions to the reals, pixel 0 as the evidence position and pixel 1 as the context position --- not to simulate a real task but to make every algebraic step of the argument chain machine-checkable.

The verification follows a completeness principle: every premise not explicitly given in the paper is explicitly included as a premise parameter of the theorem, and nothing is implicitly assumed. This principle directly shapes the verified form of Lemma 3, detailed below. The guaranteed scope of the artifact is exactly the guarantee-tier table of Section~\ref{sec:layers}. Universal claims beyond that scope, for example the satisfaction of Criterion P4 on real tasks, are not covered by the machine verification. The sites where the verified premises differ from the prose are stated explicitly in the next section: Criterion P4 enters the development as the parameterized premise $g(R) > \epsilon$ of Lemma 3 rather than as a formalized criterion; the component form of Section~\ref{sec:p4}, from which the risk-gap floor is derived, is prose that the artifact does not formalize; and the per-position coefficient \eqref{eq:focgeneral} of a general pairwise potential is derived on the toy carrier, while its continuous-side counterpart, the passage from the variational objective to pointwise form, is a correspondence of calculus of variations rather than a machine-checked step.

\subsection{Three Levels of Verification Strength}

The verified content divides into three levels of strength.

The strongest level is complete deduction, with conclusions and proofs free of extra premises. It includes the strengthened form of Lemma~1, where a single global scalar at an arbitrary fitting point forces residuals and consistency terms to be proportional across positions, \eqref{eq:proportion}, the proof being two substitutions and one commutativity; the two-sided form of Lemma~2, where Paradigm~E with a frozen scalar coefficient has no nonzero stationary point at a fitting point under heterogeneity while the objective of Paradigm~I always has a unique minimizer; the classification of Proposition~4, whose proof unfolds line by line along the collision argument of Proposition~4; the derivation of the general first-order condition for an arbitrary even differentiable potential; the demand-ratio algebra of Lemma~1's general form on the indexed family; and the closure of Theorem~5.

The next level is the landing of the criteria. Ill-posedness is formalized as the observation mask having a hole, and both directions of the mask criterion are machine-provable: a hole implies ill-posedness by constructive bifurcation, and a full mask implies unique determination by position-wise pinning. The mask form of Criterion P3 constructs the collision sample pair with evidence pinned and context freely bifurcating: take two structure fields that take the observed value at evidence positions and two different numbers at context positions; both fields agree with the observation on the mask. That P3 implies P1 is thereby also provable. The collision pair feeds Theorem 5, and the heterogeneity Criterion P2 is consumed at Lemmas 1 and 2.

The weakest level is premise parameterization, which falls on Lemma 3. The toy form lists upper-tail realization, a risk gap exceeding the bound, and nonnegativity of the minimum risk all explicitly as theorem premises; in the general measure form, upper-tail realization is given by Mathlib's first-moment-method theorem and is no longer a premise, leaving only a risk gap exceeding the bound and nonnegativity of the minimum risk as explicit premises. The verification side parameterizes the risk-gap premise as $g(R) > \epsilon$, which is the fixed-tolerance instance of the risk-gap floor derived in Section~\ref{sec:p4}: the machine consumes the positivity of the gap at the tolerance under discussion, while the derivation of that positivity from the barrier, bounded-description, and realization conditions is the prose component recorded as Tier~C in Section~\ref{sec:layers}. The conclusion is that the generalization bound cannot hold, and the proof is exactly the four-line derivation of Lemma 3. Among these, the risk gap exceeding the bound is demonstrated in mechanism by the toy instance, whose numbers Example~\ref{ex:gap} works out in full: the frozen family cannot escape a positive class gap on the three-point sample, and the evidence-computable half of that example, together with the constructive rule that attains the optimal risk, is machine-checked alongside it.

\subsection{The Final Theorem and the Scope Statement}

The final theorem \texttt{finale} composes the two sides into one conjunction: under Criterion P3 and the measure-theoretic premises of Lemma 3, a rule carried by fixed parameters neither belongs to Mode A nor can satisfy the generalization bound on the open set. The carrier side answers ``where the rule comes from'', and the generalization side answers ``what goes wrong with the rule on unseen inputs''; together they are the algebraic form of the negation of ``solo use is complete''. The guaranteed scope of the verification is given component by component in the next section.

\subsection{Correspondence from the Toy Carrier to the Continuous Case}
\label{sec:layers}

A blanket ``from toy to real'' declaration leaves the reader with a question of trust; this section upgrades it to a component-wise guarantee. Every argument component is assigned by guarantee strength to one of three tiers. Tier A is the carrier-independent algebraic core, whose proofs cite no property of the carrier and can be machine-proved directly at general types; the toy end and the continuous end are both mere instantiations. Tier B is structured general forms, which require the structure of measure spaces, inner product spaces, or modules, for which Mathlib already provides anchors, and which are machine-provable up to the level of general structures. Tier C is the remaining modeling components, belonging to derivations or correspondences on the continuous side, not covered by the machine guarantee and listed explicitly item by item. Table~\ref{tab:layers} records the tier of each component.

\begin{table}[t]
\caption{Guarantee tiers of the argument components. Tier A: carrier-independent algebraic core, proved at general types. Tier B: structured general forms anchored on Mathlib structure theorems. Tier C: remaining modeling components, not covered by the machine guarantee.}
\label{tab:layers}
\centering
\small
\begin{tabular}{@{}>{\raggedright\arraybackslash}p{0.17\textwidth}>{\raggedright\arraybackslash}p{0.19\textwidth}>{\raggedright\arraybackslash}p{0.40\textwidth}>{\centering\arraybackslash}p{0.12\textwidth}@{}}
\hline
\textbf{Component} & \textbf{Toy form} & \textbf{General or continuous form} & \textbf{Tier}\\
\hline
Ill-posedness criterion & Four-position mask with a hole & Necessary-and-sufficient criterion via nontrivial kernel of a linear observation operator & B\\
\hline
Collision construction & Four-position mask bifurcation & Pointwise bifurcation over arbitrary index types; bump construction on the continuous side & A and B\\
\hline
Algebraic core of Lemma 1 & Proportionality over consistency terms & Proportionality identity for arbitrary real quintuples & A\\
\hline
Endogenous side of Lemma 2 & Unique minimizer of a sum of squares & Positive definiteness of inner-product-space norms, including the almost-everywhere form for square-integrable functions & B\\
\hline
Measure framework of Lemma 3 & Three-point toy samples & Arbitrary probability measures; upper-tail realization given by the first-moment-method theorem & B\\
\hline
Classification and closure & Nondegenerate instance of the linear form & Nondegenerate fixed rules over arbitrary index types do not belong to Mode A & A\\
\hline
Objective to first-order condition & Coordinate derivative of the pairwise term, and Fermat's theorem along a coordinate line & Per-position coefficient of an arbitrary even differentiable potential; the paper's normalization of constant factors recovered exactly & A\\
\hline
Context rule family (Definition~2) & Indexed family over the four positions, with its frozen and constant specializations & Family indexed by position, carried with the general stationarity algebra of the demand ratio; the constant family is the specialization Lemma~1 excludes & A\\
\hline
Continuous-side optimality derivation & Not applicable & Euler--Lagrange equation and the operator correspondence of the consistency term & C\\
\hline
Network nondegeneracy & Nonzero weights & Checkable by evaluating on collision inputs & C\\
\hline
Component form of Criterion P4 & The gap premise is a parameter of Lemma 3 & Barrier, bounded-description, and realization conditions, and the barrier-to-gap inequality of Section~\ref{sec:p4} & C\\
\hline
\end{tabular}
\end{table}

Three remarks on the table. First, Tier~A is a mathematical fact rather than a defence: the proof of Lemma~1 is two substitutions and a commutativity, and its content is a proportionality identity over arbitrary reals; classification and closure likewise consume only that values differ on collision pairs, over arbitrary index types. Second, Tier~B pushes the distributional components onto Mathlib structure theorems: upper-tail realization is the first moment method, so it is demoted from a premise of Lemma~3 to the theorem \texttt{tail\ttu realization}, and the premises of the general form of Lemma~3 reduce to a risk gap exceeding the bound and nonnegativity of the minimum risk; the unique minimizer of Paradigm~I comes from positive definiteness of inner product spaces, with almost-everywhere equality in the measure form. That reduction from three premises to two is the one step where the machine crosses the toy-to-general boundary. Third, Tier~C has three items: the continuous-side derivation of first-order optimality conditions, whose discrete counterpart, the per-position coefficient \eqref{eq:focgeneral} of a general pairwise potential, is machine-checked, and whose boundary is the operator-valued coupling that the general form of Lemma~1 does not claim; network nondegeneracy on collision inputs, which this paper neither assumes nor machine-proves; and the component form of Criterion P4, namely the barrier, bounded-description and realization conditions together with the barrier-to-gap inequality, which the artifact does not formalize because it consumes the gap premise directly as the parameter of Lemma~3.

The scope statement is thereby precise to the component. Machine verification guarantees all of Tier~A and the structured forms of Tier~B: the algebraic core is machine-proved at general types, and the measure and inner-product-space parts rest on Mathlib theorems. The three Tier~C components are not covered: the continuous-side variational derivation is a standard citation and its discrete counterpart, the per-position coefficient \eqref{eq:focgeneral} with the first-order condition it yields, is the machine-checked part, while the operator-valued couplings outside that form are not claimed; network nondegeneracy is a checkable property; and the component form of Criterion~P4 with the inequality that derives the floor is carried by prose.

\subsection{Artifact Availability}

The project repository is \texttt{https://github.com/\hspace{0pt}aki-xavier/\hspace{0pt}frozen-axis}, and the Lean artifact is published there in the \texttt{lean-proof} directory: a self-contained Lake project whose \texttt{lakefile.lean} and \texttt{lake-manifest.json} pin Mathlib and every transitive dependency to an exact revision, and whose \texttt{lean-toolchain} pins the toolchain to \texttt{leanprover/\hspace{0pt}lean4:v4.34.0}; the reproduction command is \texttt{lake build}, run from inside that directory. The development file is 1219 lines and contains all definitions, lemmas, propositions, the final theorem, the example variants and the generalization upgrades, and its content is identified by the blob hash \texttt{5c6cc73a}, so the verified object is fixed independently of any later revision of the repository.

The build carries three properties worth stating rather than leaving to inspection: it completes with zero errors and zero warnings; the development contains no \texttt{sorry} and no \texttt{axiom} declaration of its own; and, strongest, \texttt{\#print axioms} on the load-bearing declarations --- \texttt{finale}, \texttt{lemma1\_sharp}, \texttt{tail\_realization} --- returns \texttt{propext}, \texttt{Classical.choice} and \texttt{Quot.sound}, Lean's standard set, with no \texttt{sorryAx}, so no step of the chain rests on an unproved assertion. The artifact is therefore auditable at the level of the proof term rather than only at that of a reported build outcome.

The artifact index is published as an online appendix accompanying this paper, so that the verification is inspectable without access to the artifact: it records the layout, the reproduction command and the axiom budget, it maps every paper item --- Definitions 1 to 4, the mask criterion, Criteria P1 to P4, Lemmas 1 to 3, Proposition 4 and Theorem 5, and the final theorem --- to the Lean declaration that discharges it, with the guarantee tier of Table~\ref{tab:layers}, and it prints the machine-checked statements of the load-bearing declarations with their premises in full, so that the premise list of \texttt{finale} can be compared line by line against the prose of Section~\ref{ch:falsify}.
